\begin{abstract}
When a language model continues a chain of reasoning that already contains a false step, it can catch the error or build on it. We study this with an audited perturb-and-continue protocol: a model solves a short program step by step, we replace one line of its own trace with a false value, let it continue with no instruction to check, and grade the continuation by re-executing the program. Whether the model rejects the false step or reuses it is decided by one property of the error. A false value that contradicts a number written earlier in the trace is rejected by capable models; a false value that could only be exposed by redoing a computation is reused. Across thirteen models from five developers, from a 1.5-billion-parameter open-weight model to frontier systems, computed falsehoods are absorbed in 88 to 100 percent of continuations by every model that reliably solves the base task, 97 to 100 percent for ten of the eleven, while the strongest models reject readable ones in 99 percent; the two weakest models solve too few base tasks to establish the contrast. Reading checks improve sharply with model quality; recomputation checks do not improve at all. The boundary is a step function: for the frontier model, absorption jumps from 2 percent to 100 percent as soon as one operation separates the model from the truth. Showing the full derivation next to the false value does not restore checking, so the failure is deference to computed-looking content rather than an inability to recompute. Four escalating instructions to verify, including a worked example, leave absorption at 100 percent. Enabling test-time reasoning drives readable-error absorption to the floor and dents only the deepest computed cells; a one-operation computed error still passes at 99 percent. Agent pipelines and chain-of-thought monitors assume errors have some chance of being caught downstream; for committed computed values---values that stand as completed steps---that chance is zero.
\end{abstract}
